\documentclass[a4paper,11pt]{article}
\usepackage[utf8]{inputenc}
\usepackage[T1]{fontenc}
\usepackage[french]{babel}
\usepackage{lmodern}
\usepackage{amsmath,amssymb,amsthm}
\usepackage{mathrsfs}
\usepackage{enumitem}
\usepackage{multicol}

\setlist[enumerate]{label=\textbf{\textcolor{Couleur8}{\arabic*.}}, left=1em}

% Si vous voulez aussi modifier les listes enumerate imbriquées :
\setlist[enumerate,2]{label=\textcolor{Couleur8}{\alph*)}}
\setlist[enumerate,3]{label=\textcolor{Couleur8}{\roman*)}}
\setlist[enumerate,4]{label=\textcolor{Couleur8}{\Alph*)}}
\usepackage{graphicx}
\usepackage{xcolor}
\usepackage{tcolorbox}
\usepackage{geometry}
\usepackage{fancyhdr}
\usepackage{titlesec}
\usepackage{cancel}
\usepackage{ifthen}
\usepackage{tikz}
\usetikzlibrary{arrows.meta}
\usepackage{tkz-tab}
\usepackage{eurosym}

% OPTION PROF/ÉLÈVE
% Décommentez UNE des deux lignes suivantes :
\newboolean{prof}\setboolean{prof}{true}   % Version professeur (avec corrections des devoirs)
%\newboolean{prof}\setboolean{prof}{false}  % Version élève (sans corrections des devoirs)

\definecolor{Couleur1}{HTML}{4A5D7A} %Th
\definecolor{Couleur2}{HTML}{A5B5D6} %Prop
\definecolor{Couleur3}{HTML}{A8C68A} %Méthode
\definecolor{Couleur6}{HTML}{8A9CC1} %Def
\definecolor{Couleur7}{HTML}{E6B459} %Rq Voc
\definecolor{Couleur8}{HTML}{D36740} %Titres
\definecolor{correction}{HTML}{D36740} % Couleur pour les corrections
\definecolor{coopmathscorr}{HTML}{F15929} % Couleur corrections coopmaths

% Configuration des marges
\geometry{
	top=2cm,
	bottom=2cm,
	inner=1.5cm,
	outer=1.5cm,
}

% Police
\renewcommand{\familydefault}{\sfdefault}

% Compteurs
\newcounter{seancenum}
\newcounter{seancenumD}
\setcounter{seancenum}{1}
\setcounter{seancenumD}{1}

% Configuration des en-têtes et pieds de page
\pagestyle{fancy}
\fancyhf{}
\ifthenelse{\boolean{prof}}{
    \fancyhead[L]{\textcolor{Couleur8}{\textbf{Corrections Automatismes - Classe de seconde - Version Professeur}}}
}{
    \fancyhead[L]{\textcolor{Couleur8}{\textbf{Corrections Automatismes - Séance 21 - Classe de seconde}}}
}
\fancyhead[R]{\textcolor{Couleur8}{\textbf{2025-2026}}}
\fancyfoot[L]{\textcolor{Couleur8}{Mme Bertail et Mme Pinto}}
\fancyfoot[R]{\textcolor{Couleur8}{\thepage}}
\renewcommand{\headrulewidth}{1pt}
\renewcommand{\headrule}{\hbox to\headwidth{\color{Couleur8}\leaders\hrule height \headrulewidth\hfill}}

% Environnements personnalisés
\newtcolorbox{seance}[1]{
    colback=Couleur6!10,
    colframe=Couleur1!65,
    fonttitle=\bfseries\large,
    title=Correction Séance \theseancenum #1,
    left=5mm,
    right=5mm,
    top=3mm,
    bottom=3mm,
    arc=0mm,
    after=\stepcounter{seancenum}
}

\newtcolorbox{seanceD}[1]{
    colback=Couleur6!5,
    colframe=Couleur6!45,
    fonttitle=\bfseries\large,
    title=Correction Devoirs - Séance \theseancenumD #1,
    left=5mm,
    right=5mm,
    top=3mm,
    bottom=3mm,
    arc=0mm,
    after=\stepcounter{seancenumD}
}

\begin{document}

\setcounter{seancenum}{21}
\newpage
\begin{seance}{} %21

\begin{enumerate}
\item  Développons l'expression :\\
  $\begin{aligned}A &=\left(-\dfrac{4}{5}c+4\right)\left(-\dfrac{6}{7}c+2\right)-\left(-\dfrac{2}{7}c-5\right)^2\\&=\left({\color{red}\boldsymbol{\dfrac{-6c}{7}}}\times {\color{red}\boldsymbol{\dfrac{-4c}{5}}}
    +{\color{blue}\boldsymbol{\dfrac{-6c}{7}}}\times {\color{blue}\boldsymbol{4}}
    +{\color{blue}\boldsymbol{2}}\times {\color{blue}\boldsymbol{\dfrac{-4c}{5}}}
    +{\color[HTML]{008002}\boldsymbol{2}}\times {\color[HTML]{008002}\boldsymbol{4}}\right)-\left({\color{red}\boldsymbol{\left( \dfrac{-2c}{7}\right) ^2}}+{\color{blue}\boldsymbol{2\times \left( \dfrac{-2c}{7}\right) \times \left( -5\right) }}+{\color[HTML]{008002}\boldsymbol{\left( -5\right) ^2}}\right)\\&={\color{red}\boldsymbol{\frac{24c^2}{35}}}{\color{blue}\boldsymbol{+\frac{-24c}{7}}}{\color{blue}\boldsymbol{+\frac{-8c}{5}}}{\color[HTML]{008002}\boldsymbol{+8}}{\color{red}\boldsymbol{+\frac{-4c^2}{49}}}{\color{blue}\boldsymbol{+\frac{-20c}{7}}}{\color[HTML]{008002}\boldsymbol{-25}}\\&=({\color{red}\boldsymbol{\frac{24c^2}{35}+\frac{-4c^2}{49}}})+({\color{blue}\boldsymbol{\frac{-24c}{7}+\frac{-8c}{5}+\frac{-20c}{7}}})+({\color[HTML]{008002}\boldsymbol{8-25}})\\&={\color[HTML]{f15929}\boldsymbol{\dfrac{148}{245}c^2-\dfrac{276}{35}c-17}}\end{aligned}$

\item  ${\color[HTML]{f15929}\boldsymbol{10^{1}}} \leqslant 84{,}074 \leqslant {\color[HTML]{f15929}\boldsymbol{10^{2}}}$ car $10^{1} = 10$ et $10^{2} = 100.$

\item  $\overrightarrow{u}=\overrightarrow{AB}$ $\Leftrightarrow\begin{cases}1=x_B-6\\-3=y_B-(-4)\end{cases}$ $\Leftrightarrow\begin{cases}1+6=x_B\\-3-4=y_B\end{cases}$ $\Leftrightarrow\begin{cases}x_B=7\\y_B=-7\end{cases}$, soit $B({\color[HTML]{f15929}\boldsymbol{7}}\,;\,{\color[HTML]{f15929}\boldsymbol{-7}})$.

\item  Comme $f(15)=60$, le coefficient $a$ tel que de $f(x)=ax$ est :\\$a=\dfrac{60}{15}=4$.\\Donc $f(x)=$ ${\color[HTML]{f15929}\boldsymbol{4x.}}$

\item \begin{enumerate}[itemsep=1em]
	\item L'équation est de la forme $x^3=k$ avec $k=125$. \\
              Le nombre dont le cube est $125$ est $5$ car $5^3=125$.\\
              Ainsi,   $S=\{5\}$.
              
	\item L'équation est de la forme $x^2=k$ avec $k=16$. Comme  $16>0$ alors l'équation admet deux solutions : $-\sqrt{16}$ et $\sqrt{16}$.\\
                Comme $-\sqrt{16}=-4$ et $\sqrt{16}=4$ alors
                les solutions de l'équation peuvent s'écrire plus simplement : $-4$ et $4$.\\
                Ainsi,  $S={\color[HTML]{f15929}\boldsymbol{\{-4\,;\,4\}}}$.
\end{enumerate}

\end{enumerate}

\end{seance}

\end{document}